\begin{frame}
\frametitle{Kódování instrukcí procesoru MIPS}

Starší typ procesoru, má tři typy instrukcí:

\begin{tabular}{|c|l|l|l|l|l|l|}\hline
Typ & 31 ... 26 & 25 ... 21 & 20 ... 16 & 15 ... 11 & 10 ... 6 & 5 ... 0 \\ \hline
R & opcode(6) & rs(5) & rt(5) & rd(5) & shamt(5) & funct(6) \\ \hline
I & opcode(6) & rs(5) & rt(5) & \multicolumn{3}{c|}{ immediate(16)} \\ \hline
J & opcode(6) & \multicolumn{5}{c|}{ address(26)} \\ \hline
\end{tabular}

\begin{itemize}
\item typ instrukce (R,I,J) je určen 6-bity opcode
\item rs -- register zdroje -- source; rd -- register cíle -- destination; rt -- někdy zdroj, někdy cíl
\item immediate, address -- konstanty použité pro operaci, immediate pro počítání, address pro skoky 
\item shamt -- konstanta využitá pouze pro posuny (shift, operace \texttt{<<}, \texttt{>>})
\item funct -- specifikace prováděné operace, např. sčítání, odčítání, posuny, atd.
\end{itemize}

\end{frame}
